
\documentclass[conference]{IEEEtran}
\IEEEoverridecommandlockouts

% IEEE Conference Template - June 2024

\usepackage{cite}
\usepackage{amsmath,amssymb,amsfonts}
\usepackage{graphicx}
\usepackage{textcomp}
\usepackage{xcolor}
\usepackage{url}

\begin{document}
	
	\title{Analyzing Class Dependencies in Java Bytecode
		for Graph-Based Software Representation}
	
	\author{
		\IEEEauthorblockN{Christian-Alexandru Dragoi}
		\IEEEauthorblockA{
			\textit{Faculty of Automation and Computers}\\
			\textit{Politehnica University Timisoara}\\
			Timisoara, Romania\\
			christian.dragoi@student.upt.ro
		}
	}
	
	\maketitle
	
	\begin{abstract}
		Understanding the structure of large object-oriented
		software systems is an important task in software
		maintenance and program comprehension. One approach
		to addressing this problem is representing software
		systems as dependency graphs, where classes are
		modeled as nodes and their relationships as edges.
		Such representations can support architectural analysis
		and the identification of key classes.
		
		This paper outlines a study of class dependency
		extraction from compiled Java bytecode. The study
		focuses on identifying which dependency types can
		be recovered from Java Archive (JAR) files and
		examining existing tools for bytecode analysis.
		A research question and an evaluation methodology
		are proposed to guide the development of a
		graph-based representation of Java software systems.
	\end{abstract}
	
	\begin{IEEEkeywords}
		Java bytecode, static analysis, class dependencies,
		dependency graphs, program comprehension,
		software architecture
	\end{IEEEkeywords}
	
	
	\section{Introduction}
	
	Large object-oriented software systems often contain
	hundreds or thousands of interconnected classes.
	Understanding the relationships between these
	classes can be difficult, especially when documentation
	is incomplete or does not reflect the current
	implementation.
	
	Program comprehension and software architecture
	reconstruction attempt to address these difficulties
	by extracting structural information directly
	from software systems. One possible approach
	is to represent a system as a dependency graph,
	where nodes correspond to classes or interfaces
	and edges describe dependencies between them.
	
	Such graphs can subsequently be analyzed using
	structural metrics and graph-based algorithms
	to identify important components of a system,
	commonly referred to as key classes \cite{b1}.
	
	Existing approaches frequently rely on source
	code analysis to extract these relationships.
	However, Java applications are commonly
	distributed as compiled JAR files, which
	contain Java bytecode rather than source code.
	
	This creates an opportunity to investigate
	which structural dependencies can be recovered
	directly from compiled artifacts.
	
	The objective of this paper is to establish
	a theoretical foundation for extracting
	class-level dependencies from Java bytecode
	and representing them as a graph.
	
	
	\section{State of the Art}
	
	\subsection{Key Class Identification}
	
	The identification of key classes is an
	important research direction in software
	maintenance and program comprehension.
	
	Sora and Chirila \cite{b1} proposed an approach
	for identifying key classes using static
	analysis and graph-based metrics. Their work
	represents software systems as graphs and
	investigates different measures of class
	importance, including direct connections
	and network centrality.
	
	Their experimental evaluation indicates that
	weighted connection metrics and adapted
	PageRank measures can be effective for
	identifying structurally important classes.
	
	A different approach was presented by
	Vale and Maia \cite{b2}, who proposed
	Keecle, a technique based on dynamic
	analysis and method invocation traces.
	
	While static analysis examines the structure
	of software without executing it, dynamic
	analysis considers relationships observed
	during execution.
	
	Both approaches demonstrate the importance
	of class relationships for understanding
	software architecture.
	
	\subsection{Class Dependency Types}
	
	Sora and Chirila \cite{b1} distinguish
	several categories of static dependencies
	between object-oriented classes.
	
	These include inheritance, interface
	implementation, field declarations,
	method parameters, return types,
	local variables, field accesses,
	method invocations, object instantiation,
	type casting and generic type bindings.
	
	A class may depend on another class
	through multiple mechanisms simultaneously.
	For example, a class may declare a field
	of another type while also invoking
	methods belonging to that type.
	
	Consequently, a dependency graph may
	need to preserve both the direction
	and the type of each relationship.
	
	Furthermore, different dependency types
	may have different structural significance.
	The use of weighted edges provides a
	possible mechanism for representing
	these differences \cite{b1}.
	
	\subsection{Java Bytecode Analysis}
	
	Java source code is compiled into class
	files containing information defined
	by the Java Virtual Machine (JVM)
	specification \cite{b3}.
	
	Class files contain structural metadata,
	including class names, superclass
	information, implemented interfaces,
	fields and method descriptors.
	
	They also contain a constant pool
	and may contain bytecode instructions
	describing method invocations,
	field accesses and object creation.
	
	However, not all source-level information
	is necessarily preserved after compilation.
	
	For example, local variable metadata
	may be absent when the corresponding
	optional debug attributes are not included.
	Additionally, dynamically resolved
	relationships, including certain uses
	of reflection, may not be reliably
	identified through static analysis alone.
	
	These characteristics must be considered
	when extracting dependency information
	from compiled Java applications.
	
	\subsection{Existing Bytecode Analysis Tools}
	
	Several tools and libraries provide
	mechanisms for examining Java bytecode.
	
	ASM \cite{b4} is a bytecode manipulation
	and analysis framework that allows
	programmatic access to class structures
	and bytecode instructions.
	
	Apache Commons BCEL \cite{b5} provides
	an alternative API for reading, analyzing
	and modifying Java class files.
	
	The JDK Class-File API \cite{b6} is
	another relevant solution, providing
	a standard API for parsing and examining
	class files.
	
	Additionally, the JDK includes tools
	such as \texttt{jdeps}, which identifies
	dependencies between classes and packages,
	and \texttt{javap}, which disassembles
	compiled class files \cite{b7,b8}.
	
	These solutions offer different levels
	of abstraction and functionality.
	
	Their suitability for constructing a
	detailed class dependency graph requires
	further investigation.
	
	
	\section{Research Question}
	
	The study aims to investigate the
	feasibility of reconstructing class-level
	dependencies directly from compiled Java
	applications.
	
	The main research question is:
	
	\begin{quote}
		\textit{RQ1: Which types of class-level
			dependencies can be reliably extracted
			from Java bytecode, and what limitations
			affect their representation in
			dependency graphs?}
	\end{quote}
	
	To address this question, the study
	will examine the dependency categories
	identified in previous research and
	the information available in Java
	class files.
	
	A secondary objective is to identify
	which existing bytecode analysis tools
	are suitable for supporting this
	extraction process.
	
	
	\section{Proposed Methodology}
	
	The proposed study will combine
	documentation analysis with a small
	experimental evaluation.
	
	\subsection{Dependency Classification}
	
	The first stage will examine the
	dependency types presented in
	the literature \cite{b1}.
	
	Each dependency category will be
	mapped to the corresponding information
	available in the JVM class file format.
	
	The dependencies will be classified
	according to whether they can be
	extracted directly, require additional
	analysis, or cannot be reliably recovered
	using static bytecode analysis alone.
	
	\subsection{Tool Analysis}
	
	The second stage will investigate
	existing Java bytecode analysis
	libraries, particularly ASM, BCEL
	and the JDK Class-File API.
	
	The tools will be compared based on
	their ability to access class metadata,
	method signatures, bytecode instructions
	and dependency-related information.
	
	Additional criteria will include
	API complexity and suitability for
	integration into a dependency
	extraction application.
	
	\subsection{Planned Experimental Evaluation}
	
	A small Java application containing
	representative dependency relationships
	will be developed and compiled
	into a JAR file.
	
	The application will contain examples
	of inheritance, interface implementation,
	field dependencies and method calls.
	
	The compiled classes will be examined
	using existing bytecode analysis tools.
	
	The extracted information will then
	be compared with the dependencies
	explicitly defined in the source code.
	
	This comparison will help determine
	which dependency categories can be
	recovered and which information may
	be lost or require additional analysis.
	
	
	\section{Results and Discussion}
	
	% TODO:
	% This section will be completed after
	% performing the experimental evaluation.
	%
	% Planned content:
	% - Dependency extraction results
	% - Comparison of bytecode analysis tools
	% - Dependency coverage and limitations
	% - Discussion of implications for graph modeling
	
	
	\section{Conclusions and Future Work}
	
	% TODO:
	% Complete this section after the study.
	%
	% Planned content:
	% - Answer to the research question
	% - Summary of supported dependency types
	% - Limitations of bytecode-based analysis
	% - Selection of a suitable parsing library
	% - Future integration with graph-based
	%   key class identification methods
	
	
	\begin{thebibliography}{00}
		
		\bibitem{b1}
		I. Sora and C.-B. Chirila,
		``Finding key classes in object-oriented software
		systems by techniques based on static analysis,''
		\textit{Information and Software Technology},
		vol. 116, Art. no. 106176, 2019.
		doi: 10.1016/j.infsof.2019.106176.
		
		\bibitem{b2}
		L. do Nascimento Vale and M. de Almeida Maia,
		``Key classes in object-oriented systems:
		Detection and assessment,''
		\textit{International Journal of Software
			Engineering and Knowledge Engineering},
		vol. 29, no. 10, pp. 1439--1463, 2019.
		doi: 10.1142/S0218194019500451.
		
		\bibitem{b3}
		T. Lindholm, F. Yellin, G. Bracha,
		A. Buckley, and D. Smith,
		\textit{The Java Virtual Machine Specification},
		Java SE 26 Edition, Oracle, 2026.
		[Online]. Available:
		\url{https://docs.oracle.com/javase/specs/jvms/se26/html/}
		
		\bibitem{b4}
		OW2 Consortium,
		``ASM: Java bytecode manipulation
		and analysis framework.''
		[Online]. Available:
		\url{https://asm.ow2.io/}
		
		\bibitem{b5}
		Apache Software Foundation,
		``Apache Commons BCEL.''
		[Online]. Available:
		\url{https://commons.apache.org/proper/commons-bcel/}
		
		\bibitem{b6}
		Oracle,
		``Package java.lang.classfile,''
		\textit{Java SE 26 API Specification}.
		[Online]. Available:
		\url{https://docs.oracle.com/en/java/javase/26/docs/api/java.base/java/lang/classfile/package-summary.html}
		
		\bibitem{b7}
		Oracle,
		``The jdeps command,''
		\textit{JDK 26 Tool Specifications}.
		[Online]. Available:
		\url{https://docs.oracle.com/en/java/javase/26/docs/specs/man/jdeps.html}
		
		\bibitem{b8}
		Oracle,
		``The javap command,''
		\textit{JDK 26 Tool Specifications}.
		[Online]. Available:
		\url{https://docs.oracle.com/en/java/javase/26/docs/specs/man/javap.html}
		
	\end{thebibliography}
	
\end{document}
